\chapter{Continuity of the Person}
\label{ch:28}

Suppose that the medical program of the preceding chapter, or the reliability program of the next, were to succeed to an extent that now seems fanciful, and that a human being could continue for ten thousand years or a million. A question that has been put off would then become practical: in what sense would the person at the end be the person at the beginning? The question has a long history in philosophy, where it is called the problem of personal identity, and it is not an idle one for a program that proposes to preserve persons. A program to preserve something must say what the something is. If the answer is the body, then the program is a program of biological repair. If the answer is the pattern of memory, character, and intention, then it may be a program that could in principle be carried out by other means, such as the replacement of tissue by engineered substitutes, or the transfer of the pattern to another substrate, and it raises the possibilities of copying, branching, and loss that the biological program does not. The answer given in this chapter is not a metaphysical one. It identifies the features of continuity that the institutions of the program rely upon, shows how those features behave over very long intervals, and states the demands that follow for the arrangements of responsibility and consent.

The classical positions are two. On the view associated with Locke, a person persists as long as there is continuity of consciousness, which later writers have generalized to psychological continuity: the connections of memory, intention, belief, desire, and character that link a person at one time with the person at another. On the competing view the person is an organism, and persists as long as the organism does, regardless of what happens to the psychology. The controversy is old and unresolved, and the cases that pull the two views apart, such as severe dementia, or the thought experiments of a teleporter that destroys a body and rebuilds it elsewhere, are cases in which intuitions differ. The most influential modern treatment, that of Parfit in 1984~\cite{parfit1984}, argued that what matters in survival is not identity but a relation, which he called Relation R, of psychological connectedness and continuity with the right kind of cause. Connectedness is the holding of direct psychological connections, such as a memory of an experience or the carrying out of an earlier intention; continuity is the holding of overlapping chains of strong connectedness. Connectedness comes in degrees and diminishes with time. I am now strongly connected with myself of last week, weakly with myself of twenty years ago, and not at all with the infant who bears my name. Yet we speak of a single life, because the chain of overlapping connections is unbroken, and Parfit's observation is that it is the chain, not the end-to-end connection, that carries what we care about.

The relevance of this account to the program is that it is a model that can be made quantitative, and the quantities it yields are instructive. Take the simplest assumption, that the strength of connection between two stages of a person separated by a time $\Delta$ decays exponentially with a characteristic time $\tau$, and that two stages are strongly connected if the strength exceeds one half. Then each stage is strongly connected with those within a window of $\tau\ln2$ on either side, and the stages form a chain of overlapping windows that extends indefinitely, whatever the total length of the life. For $\tau=20$ years the window is $13.9$ years. A person who lives a million years consists of about $72{,}000$ such windows, and is continuous with its earliest stage by a chain of $72{,}000$ links, but the direct connectedness between the first and the last is $e^{-50{,}000}$, a number with no practical meaning, and it falls below one part in a thousand after only $138$ years. The structure is that of a rope in which no fiber runs the full length. Parfit would say that the person is the same in the sense that matters, and a critic would say that the person is a succession of different people each of whom is connected to the next. The program need not decide the question. It needs only to observe that the practices of accountability, promise, and consent rest on the connectedness of stages within a window, and that it is these practices, and not the existence of a chain, that long lives test.

The memory that carries connectedness is itself finite. If a person acquires information at a rate $r$ and can retain a total amount $M$, then the whole of the record is retained only for a time $M/r$, after which the person must discard or compress. For stipulated and merely illustrative values of $M=10^{15}$ bits and $r=10^{10}$ bits per year, the interval is $10^5$ years, and beyond it the fraction of life retained falls inversely with the length of the life. The values are not established, since estimates of the capacity of the brain differ by several orders of magnitude, and the point of the calculation is the form, not the figure: an unaugmented brain would retain an ever-smaller fraction of a long life, and any augmentation of memory, whether by external records or by modification, would become part of what the person is. The practical consequences are those of the philosophical literature on the extended mind, applied at scale. A person of great age will depend on records, and the integrity, ownership, and privacy of those records become matters of the highest importance, since their alteration would alter the person.

Three consequences for the institutions of the program deserve statement. First, responsibility across long intervals requires a rule. A person who made a promise or committed a wrong five thousand years earlier is connected to the doer by a chain and not by memory or character, and the practices of law have in the human case limited responsibility over time by statutes of limitation and by the rule that the young are not bound by the commitments of their earlier years in the way that an adult is. The program does not prescribe particular limits, but it requires that rules of this kind exist and be public, and that the duty-bearer relations of Chapter 2 be defined so as to remain well posed when a bearer may persist for millennia: the coverage and reassignment conditions must hold across the life of the bearer. Second, consent given by a stage cannot be taken to bind stages beyond its window of connectedness without renewal. An agreement entered into by a person at one age, with consequences that extend for centuries, is an agreement with respect to a future self that may be at the outset only distantly connected, and the sensible rule is that long-term commitments require periodic reaffirmation, with a right to withdraw at stated intervals. Third, the possibility of branching, of copies of a pattern that are equally continuous with the original, destroys the one-to-one relation on which law and property depend. Relation R can hold between a person and two later persons at once, and it is not transitive in the way that identity is. The program therefore prohibits, as a condition of authorization of any technology that could produce copies, the creation of copies without a prior determination of their standing, and treats the copy of a person as a person with its own claims, to be neither the property of the original nor the original's instrument.

A final consideration concerns what the program owes to those who decline to continue. A society in which many persons live very long lives will contain persons who prefer not to, and the freedom to choose a natural lifespan, to decline interventions, and to die, is a condition of the legitimacy of the program. The right to refuse, which Chapter 13 required of every authorization, applies here with particular force, because the interest at stake is the person's own life. The institutions must not make the declining of treatment a source of disadvantage, nor make the acceptance of treatment a condition of full membership in society.

\section*{Formalization}

Let a life consist of stages indexed by times $t_0<t_1<\dots<t_N$ with spacing $\Delta=t_{i+1}-t_i$ and pairwise connectedness $c(i,j)=\exp(-|t_i-t_j|/\tau)\in(0,1]$. Fix a threshold $\theta\in(0,1)$, with $\theta=1/2$ for Parfit's "strong" connectedness.

\begin{definition}[Continuity]
Stages $i$ and $j$ are \emph{strongly connected} if $c(i,j)\ge\theta$. A \emph{continuity chain} from $0$ to $N$ is a sequence of stages $0=i_0<i_1<\dots<i_m=N$ with consecutive stages strongly connected. The life is \emph{continuous} if a chain exists.
\end{definition}

\begin{proposition}
(i) A continuity chain from $0$ to $N$ exists if and only if $\Delta\le\tau\ln(1/\theta)$. (ii) The minimum number of links in a chain is $m_{\min}=\lceil (t_N-t_0)/(\tau\ln(1/\theta))\rceil$ when $\Delta$ divides the window. (iii) The direct connectedness of the first and last stages is $c(0,N)=e^{-(t_N-t_0)/\tau}$, and falls below $\varepsilon$ when $t_N-t_0>\tau\ln(1/\varepsilon)$; it is below $\theta$ whenever $m_{\min}\ge2$.
\end{proposition}

\begin{proof}
(i) If $\Delta>\tau\ln(1/\theta)$ then no two distinct stages are strongly connected, since $c(i,j)\le e^{-\Delta/\tau}<\theta$ for $i\ne j$, and no chain exists; if $\Delta\le\tau\ln(1/\theta)$ the consecutive stages form a chain. (ii) Each link spans at most $\tau\ln(1/\theta)$ in time, so at least $(t_N-t_0)/(\tau\ln(1/\theta))$ links are needed, and the bound is attained by taking the largest permitted steps. (iii) Direct substitution; $c(0,N)<\theta$ iff $t_N-t_0>\tau\ln(1/\theta)$ iff $m_{\min}\ge2$.
\end{proof}

For $\tau=20$ years and $\theta=\tfrac12$ the window is $13.86$ years; a life of $10^6$ years requires at least $72{,}135$ links; and the end-to-end connectedness falls below $10^{-3}$ at $138$ years and is $e^{-5\times10^4}$ at $10^6$ years. The exponential law is a stipulation, and any decay law yields a finite window, so that the qualitative conclusion, that continuity is compatible with arbitrarily small end-to-end connection, depends only on the existence of a threshold window.

For memory, let the person acquire $r$ bits per unit time and retain at most $M$ bits, with uniform retention of the full record up to the capacity. The retained fraction of a life of length $T$ is $f(T)=\min\{1,M/(rT)\}$, equal to $1$ for $T\le T_\ast=M/r$ and decreasing as $1/T$ beyond. For stipulated $M=10^{15}$ bits and $r=10^{10}$ bits yr$^{-1}$, $T_\ast=10^5$ yr and $f(10^6)=0.1$. If the person supplements biological memory by an external store of capacity $M_e$, the retained fraction is $\min\{1,(M+M_e)/(rT)\}$, so that the fraction of the long life that is accessible depends on a store whose custody is, by the argument of the chapter, constitutive of the person.

For standing, let $n$ copies of a pattern be produced at a branching event, each strongly connected to the original stage and to no other. Then the relation of strong connectedness between the original stage and each descendant holds, but the relation among descendants fails, so strong connectedness is not transitive; no assignment of a single identity to the descendants is consistent with identity as an equivalence relation unless $n=1$. This is the formal reason for the requirement that standing be determined before copies are created.
